\section{Introduction}
Weak gravitational lensing of the large-scale structure (cosmic shear) is one of the key tools in cosmology for understanding the distribution of cosmic structures and deducting the cosmological model that gave rise to it. 
According to General Relativity, light is bent when it passes through gravitational perturbations sourced by cosmic structure. 
This means that the light rays from distant galaxies are ``lensed'' as they travel to us, resulting in small distortions in the observed galaxy shapes. 
This distortion encodes information of the cosmic structure, allowing us to probe (visible and invisible) cosmic structure.

The central observable in weak lensing is thus the shapes of galaxy images. 
In particular, we are interested in the galaxy shapes after it has been distorted (or lensed) by the cosmic structure but before it is distorted further by atmospheric or instrumental effects.
As this exact quantity is not observable, we can only derive it by having a good model of the atmospheric and instrumental effects, or the Point Spread Function (PSF).

The PSF refers to how a point source will look like on our detectors at a given time and location on the focal plane. 
It encodes the net effect from the atmosphere, the optics and the detectors and can be thought of as a convolution kernel (or a Green's function) that is applied to the light that comes into the telescope. 
Having a precise and accurate PSF model for each galaxy is key to the success of weak lensing science for Rubin LSST.

One of the subtleties in modeling the PSF is that the PSF is wavelength-dependent. 
In particular, Differential Chromatic Refraction \citep[DCR,][]{DMTN-017, 2012PASP..124.1113A, 2015ApJ...807..182M} and chromatic seeing are the dominating wavelength-dependent PSF effects that exist in ground-based imaging instruments such as Rubin's LSSTCam. When we look through the atmosphere at different airmass, the light rays that enter the atmosphere come in a non-perpendicular incident angle and refraction happens at the interface.
Blue light refracts more than red light, causing small distortions or displacements of the observed objects that depend on their wavelength. 
This tech note focuses on this effect and how we could model it in the Rubin pipeline.

The Rubin Observatory uses a PSF fitting strategy adapted from the Dark Energy Survey (DES), which is based on the PSF in the Full Focal plane (\texttt{PiFF}) algorithm \citep{jarvis2021}. 
In this note we examine the performance of the PSF modeling in the Rubin commissioning data. 
We use the weekly intermittent LSST Science Pipelines \citep{PSTN-019} \texttt{w\_2025\_49} data.

The structure of this note is as follows. Section 2 describes the DCR effect. 
Section 3 describes the data products used in this study. 
Section 4 discusses the fitting procedure, PSF residual definitions and the PSF fitting configuration.
Section 5 shows the results of enabling chromatic fitting on PSF residuals. 
In Section 6 we compare our results with that from DES and in Section 7 propagate the PSF residuals to expected residuals in shear measurements (we only consider additive biases due to DCR and not multiplicative biases due to chromatic seeing).

